\section{Proof-producing finite traces}
\label{sec:trace}

\subsection{Trace systems}

A trace system packages a bounded computation and the invariant needed to
justify it.  All external data---the mathematical input, requested precision,
finite tables, and supplied certificate codes---are paired into one natural
number \(u\).  Fix primitive-recursive symbols
\[
 \mathsf{good}(u),\quad \mathsf{init}(u),\quad \mathsf{bound}(u),
 \quad \mathsf{step}(i,s,u),\quad \mathsf{inv}(i,s,u),
 \quad \mathsf{accept}(s,u).
\]
The run is the primitive recursion
\begin{align*}
 \Run(0,u)&=\mathsf{init}(u),\\
 \Run(i+1,u)&=\mathsf{step}(i,\Run(i,u),u).
\end{align*}
In the repository syntax this is exactly
\code{PR.prec init step}; the recursion argument is followed by the single
parameter \(u\).

The local obligations are not metatheoretic propositions.  They are supplied
as object-language derivations of the open formulas
\begin{align}
 \mathsf{good}(u)=1&\to
   \mathsf{inv}(0,\mathsf{init}(u),u)=1,
   \tag{B}\label{eq:traceB}\\
 \mathsf{good}(u)=1&\to
   \mathsf{inv}(i,s,u)=1\to
   \mathsf{inv}(i+1,\mathsf{step}(i,s,u),u)=1,
   \tag{S}\label{eq:traceS}\\
 \mathsf{good}(u)=1&\to
   \mathsf{inv}(\mathsf{bound}(u),s,u)=1\to
   \mathsf{accept}(s,u)=1.
   \tag{F}\label{eq:traceF}
\end{align}

\begin{theorem}[Embedded finite-trace theorem]
\label{thm:trace-reflection}
Given derivations of \eqref{eq:traceB}--\eqref{eq:traceF}, one can construct a
derivation
\[
 []\vdash[\PRAAx]\forall u\;
 \bigl(\mathsf{good}(u)=1\to
 \mathsf{accept}(\Run(\mathsf{bound}(u),u),u)=1\bigr).
\]
The construction is structural in the three supplied proof terms.
\end{theorem}

\begin{proof}
Keep \(u\) free and place \(G(u):\mathsf{good}(u)=1\) in the context.  Use the
quantifier-free induction formula
\[
  I(i,u):\equiv
  \mathsf{inv}(i,\Run(i,u),u)=1.
\]

For the base, instantiate \eqref{eq:traceB} at \(u\), eliminate the implication
with the context hypothesis, and obtain the invariant at
\(\mathsf{init}(u)\).  The checked theorem \code{PRA.precZero} derives
\(\Run(0,u)=\mathsf{init}(u)\).  Equality substitution changes the state
argument and yields \(I(0,u)\).

For the successor step, work in the induction context containing \(I(i,u)\)
and the appropriately shifted copy of \(G(u)\).  Instantiate
\eqref{eq:traceS} with \(s=\Run(i,u)\) and apply both implications.  This gives
\[
 \mathsf{inv}(i+1,\mathsf{step}(i,\Run(i,u),u),u)=1.
\]
The checked theorem \code{PRA.precSucc} derives the equality between the state
inside this formula and \(\Run(i+1,u)\).  A single use of
\code{Prf.eqSubst1} yields \(I(i+1,u)\).

Apply \code{PRA.qfInd}.  Specialize the resulting universal invariant at
\(\mathsf{bound}(u)\).  Instantiate \eqref{eq:traceF} with the terminal run
state, eliminate its two implications, discharge \(G(u)\), and universally
introduce \(u\).  Every operation is either a constructor of \code{Prf} or a
checked structural theorem in \code{Deduction.lean} and \code{PR.lean}.
\end{proof}

\subsection{Reduction to the checked primitive-recursion transformer}

The repository already contains a stronger reusable component in the
\code{ipra\_local\_global} packet.  Its theorem
\code{primitiveRecursionInvariant} accepts a quantifier-free formula
\(P(s,i,\bar x)\), object proofs of a base and local step, and returns the
object proof of \(P(\Run(i,\bar x),i,\bar x)\) for all \(i\).  Therefore the
trace theorem need not reimplement binder plumbing.

Take the one-parameter formula
\[
  P(s,i,u):\equiv\mathsf{inv}(i,s,u)=1
\]
and context \([G(u)]\).  The three trace obligations produce exactly the base
and successor premises required by \code{primitiveRecursionInvariant}.  Its
specialized theorem \code{primitiveRecursionAt} is applied at
\(\mathsf{bound}(u)\); the final obligation then closes the proof.  This
reduction is important for implementation: the difficult de Bruijn
substitutions and recursion equations have already been exercised by a compiled
object-theory proof transformer.

\subsection{Certificate mode}

A large computation can be externally produced and internally verified.  A
certificate code \(c\) stores a bounded state sequence and auxiliary witnesses.
A fixed checker verifies
\begin{align*}
 s_0&=\mathsf{init}(u),\\
 s_{i+1}&=\mathsf{step}(i,s_i,u),\\
 \mathsf{inv}(i,s_i,u)&=1,\\
 \mathsf{accept}(s_B,u)&=1.
\end{align*}
All list assertions are finite folds.

\begin{theorem}[Certificate reflection]
\label{thm:certificate-reflection}
Suppose the soundness of each checker clause is supplied as an object-language
derivation.  Then a fixed derivation constructor proves
\[
 []\vdash[\PRAAx]\forall u\forall c\;
 \bigl(\mathsf{TraceOK}(u,c)=1\to
       \mathsf{Conclusion}(u,\mathsf{out}(c))\bigr).
\]
\end{theorem}

\begin{proof}
Fold correctness exposes every checked clause.  One proof route reads the final
accepted state directly and applies the supplied soundness implication.  A
second route proves by quantifier-free induction that each certified state is
equal to the canonical run and then invokes \cref{thm:trace-reflection}.  The
first route is smaller; the second is useful when later modules need equality
with the canonical run.
\end{proof}

Certificate mode separates proof complexity from producer complexity.  A
matrix factorization, isolating-disk list, high-precision numerical inequality,
finite cover, or descent trace can be produced by sophisticated software while
the iPRA proof checks only a small fixed relation.

\subsection{Composition}

A certified finite transducer consists of a trace system, an input contract, an
output term, and a derivation of its output contract.  Two transducers compose
by pairing their states and adding a phase bit.  Phase zero runs the first
system; a switch state freezes its output and any magnitude or guard metadata;
phase one runs the second.  The combined invariant is a Boolean case split on
the phase.

\begin{theorem}[Finite transducer closure]
\label{thm:transducer-closure}
Certified finite transducers are closed under sequential composition, finite
products, bounded loops, overlapping guarded branches, and finite acyclic
networks.  The composite bound, state transition, and proof are primitive
recursive in the fixed component syntax and structural in the component
proofs.
\end{theorem}

\begin{proof}
For sequential composition, use the phase construction just described.  In
phase zero, invariant preservation is the first component's step theorem.  At
the switch, its acceptance theorem establishes the second component's input
contract.  In phase one, preservation is the second component's step theorem.
Apply \cref{thm:trace-reflection}.  Products pair states and conjoin invariants.
A bounded loop is already a trace.  A guarded branch uses a total Boolean
selector and requires each overlapping branch to prove the same output
contract.  A finite acyclic network is topologically ordered and flattened
into phases.  Proof composition uses substitution, weakening, implication
elimination, and \code{Prf.trans}.
\end{proof}
